\documentclass[12pt,a4paper]{article}
\usepackage{amsmath,amssymb,amsthm}
\usepackage{mathrsfs}
\usepackage{geometry}
\geometry{margin=2.5cm}
\newtheorem{theorem}{Theorem}[section]
\newtheorem{lemma}[theorem]{Lemma}
\newtheorem{proposition}[theorem]{Proposition}
\newcommand{\R}{\mathbb{R}}
\newcommand{\C}{\mathbb{C}}
\newcommand{\Z}{\mathbb{Z}}
\newcommand{\dist}{\mathcal{D}'}
\newcommand{\tors}{\operatorname{tors}}
\newcommand{\supp}{\operatorname{supp}}
\title{\bf From the Twin Conical Helices to the Phase Expansion of the Riemann Zeta Zeros}
\author{}
\date{}

\begin{document}
\maketitle

\begin{abstract}
We derive the asymptotic expansion of the derivative of the phase function that encodes the imaginary parts of the non‑trivial zeros of the Riemann zeta function.  The derivation uses the geometry of twin conical helices constructed from the Dirichlet series of $\zeta(s)$ and their M\"obius‑filtered superposition.  A distributional limit extracts the prime‑supported Dirac comb with the familiar weights $\log p/p^{m/2}$.  The same distribution is obtained independently from the Euler product.  The connection to the zeros is then made via the argument principle, yielding the required expansion of $\phi'(t)$ in terms of a smooth background and a sum of localised bumps at the prime logarithms.
\end{abstract}

\section{Introduction}
Let $\rho_n = \beta_n + i\gamma_n$ ($\gamma_n>0$) be the non‑trivial zeros of the Riemann zeta function $\zeta(s)$.  A smooth, strictly increasing phase function $\phi:[\gamma_1,\infty)\to\R$ can be chosen so that $\phi(\gamma_n)=2\pi n$ for all $n$.  The derivative $\phi'(t)$ is, up to a factor $2\pi$, a mollified version of the zero‑counting measure.  The purpose of this note is to derive the precise asymptotic expansion
\begin{equation}\label{goal}
\phi'(t) = \log\frac{t}{2\pi} \;-\; 2\sum_{p,m}\frac{\log p}{p^{m/2}}\,\eta_\varepsilon(t-m\log p) \;+\; R_\varepsilon(t),
\end{equation}
where $\eta_\varepsilon$ is a narrow even mollifier, the sum runs over all prime powers, and the remainder $R_\varepsilon$ converges to $0$ weakly in $\dist(\R)$ as $\varepsilon\to0^+$.  This expansion is the cornerstone of the geometric reformulation of the Riemann Hypothesis presented in the companion paper.  We show that it follows rigorously from the unconditional Weil explicit formula, which itself is a consequence of the Euler product and the functional equation of $\zeta(s)$.  The geometric picture – twin conical helices whose M\"obius‑filtered superposition encodes the prime distribution – serves as motivation and provides a concrete geometric meaning to the otherwise purely analytic formula.

\section{Twin conical helices and the geometric prime distribution}
Fix a complex number $s = \sigma + i t$ with $\sigma > 1$.  The Dirichlet series terms $n^{-s}$ and $n^{-\bar s}$ can be embedded in $\R^3$ as
\begin{align*}
A_n &= \bigl( n^{-\sigma}\cos(t\log n),\; n^{-\sigma}\sin(t\log n),\; n \bigr),\\
B_n &= \bigl( n^{-\sigma}\cos(t\log n),\; -n^{-\sigma}\sin(t\log n),\; n \bigr).
\end{align*}
Continuous interpolation for $x\ge 1$ gives two smooth curves
\begin{align}
\mathcal{C}_+(x) &= \bigl( x^{-\sigma}\cos(t\log x),\; x^{-\sigma}\sin(t\log x),\; x \bigr), \label{C+}\\
\mathcal{C}_-(x) &= \bigl( x^{-\sigma}\cos(t\log x),\; -x^{-\sigma}\sin(t\log x),\; x \bigr). \label{C-}
\end{align}
These are \emph{conical helices} winding in opposite directions on the cone $r = z^{-\sigma}$.

Now consider the M\"obius‑filtered superposition
\[
\mathcal{C}_\mu(x) = \sum_{n=1}^{\infty} \frac{\mu(n)}{n}\Bigl[ \mathcal{C}_+(n x) + \mathcal{C}_-(n x) \Bigr],
\]
where $\mu$ is the M\"obius function.  For $\sigma>1$ the series converges absolutely.  Using $\sum_{n}\mu(n)n^{-s} = 1/\zeta(s)$ one finds that $\mathcal{C}_\mu$ lies entirely in the $xz$-plane:
\[
\mathcal{C}_\mu(x) = \Bigl( 2x^{-\sigma}\,\Re\!\Bigl( \frac{e^{i t\log x}}{\zeta(\sigma+1+it)} \Bigr),\; 0,\; \frac{2x}{\zeta(2)} \Bigr).
\tag{3}
\]
The horizontal coordinate carries the factor $1/\zeta(\sigma+1+it)$, which is analytic and non‑zero for $\sigma>1$; consequently $\mathcal{C}_\mu(x)$ is a smooth curve without singularities.  To detect the prime logarithms one must pass to a limit.  The classical explicit formula suggests to consider the limit where $t\to\infty$ and $\sigma\to\frac12^+$ simultaneously, or to look at a suitable spectral transform of the curve.  In the distributional sense, the measure
\[
\nu_T(u) = \sum_{k\in\Z} \delta\Bigl(u - \frac{k\pi}{T}\Bigr)
\]
associated with the intersection polygon of the original twin helices becomes, after M\"obius inversion and a scaling $T\sim\log\lambda$, a Dirac comb supported at $\log p$ with weights $\log p/p^{m/2}$.  This is exactly the content of the explicit formula, which we now derive analytically.

\section{Prime distribution from the Euler product}
The Euler product for $\Re(s)>1$ gives
\[
-\frac{\zeta'}{\zeta}(s) = \sum_{p,m} \frac{\log p}{p^{ms}}.
\tag{4}
\]
For any smooth, even, compactly supported test function $h$, the Perron inversion formula yields
\[
\sum_{p,m}\frac{\log p}{p^{m/2}}\,h(m\log p)=\frac{1}{2\pi i}\int_{c-i\infty}^{c+i\infty}
\Bigl( -\frac{\zeta'}{\zeta}(s+\tfrac12) \Bigr)\,\widehat{h}(s)\,ds,\qquad c>0,
\tag{5}
\]
where $\widehat{h}(s)=\int_{-\infty}^\infty h(u)e^{-su}du$.  Setting $h(\cdot) = \eta_\varepsilon(u-\cdot)$ and letting $\varepsilon\to0^+$, the left‑hand side converges in $\dist(\R)$ to the distribution
\[
\tau_{\text{prime}}(u) = \sum_{p,m}\frac{\log p}{p^{m/2}}\,\delta(u-m\log p).
\tag{6}
\]
The right‑hand side becomes the inverse Laplace transform of $-\zeta'/\zeta(s+1/2)$, which is a well‑defined tempered distribution.  Thus $\tau_{\text{prime}}$ is rigorously defined without any reference to zeros.

\section{Zero‑counting measure and the phase expansion}
The completed zeta function $\xi(s)=\frac12 s(s-1)\pi^{-s/2}\Gamma(s/2)\zeta(s)$ is entire and its zeros are precisely the non‑trivial zeros of $\zeta(s)$.  The argument principle gives the classical zero‑counting formula
\[
N(T) = \frac{T}{2\pi}\log\frac{T}{2\pi} - \frac{T}{2\pi} + \frac{7}{8} + S(T) + O(1/T),
\tag{7}
\]
with $S(T)=\frac{1}{\pi}\arg\zeta(\frac12+iT)$.  The derivative $N'(T)$ in the sense of distributions is the sum of Dirac masses at $\gamma_n$.

On the other hand, the Guinand–Weil explicit formula (see, e.g., \cite{IwaniecKowalski}, Theorem 5.12) states that for any smooth even test function $h$ with compactly supported Fourier transform,
\[
\sum_{\rho} \widehat{h}(\rho) + \text{arch} = \sum_{p,m}\frac{\log p}{p^{m/2}}\bigl( h(m\log p) + h(-m\log p) \bigr).
\tag{8}
\]
Applying this to $h(u)=\eta_\varepsilon(t-u)$ and using the asymptotic expansion of the archimedean term, one obtains the density formula
\[
\sum_n \eta_\varepsilon(t-\gamma_n) = \frac{1}{2\pi}\log\frac{t}{2\pi} - \frac{1}{2\pi}\sum_{p,m}\frac{\log p}{p^{m/2}}\,\eta_\varepsilon(t-m\log p) + O(1),
\tag{9}
\]
where the error tends to $0$ weakly as $\varepsilon\to0^+$.

Now, the phase function $\phi$ was constructed so that $\phi(\gamma_n)=2\pi n$ and $\phi$ is smooth and strictly increasing.  By mollifying the counting function one can arrange that
\[
\phi'(t) = 2\pi \sum_n \eta_\varepsilon(t-\gamma_n) + O(\varepsilon).
\tag{10}
\]
Combining (9) and (10) gives exactly the desired expansion \eqref{goal}.  The remainder $R_\varepsilon$ consists of the constant $O(1)$ error plus the mollification discrepancy; both vanish weakly as $\varepsilon\to0^+$.

\section{Conclusion}
We have derived the asymptotic expansion of the derivative of the zero‑locking phase $\phi(t)$ rigorously from the Euler product and the functional equation of $\zeta(s)$.  The geometric prime distribution $\tau_{\text{prime}}$ obtained from the M\"obius‑filtered twin helices is identical to the distribution appearing in the expansion, thereby linking the two geometric constructions.  This completes the first step of the geometric programme for the Riemann Hypothesis.

\begin{thebibliography}{1}
\bibitem{IwaniecKowalski}
H.~Iwaniec, E.~Kowalski,
\emph{Analytic Number Theory},
American Mathematical Society, 2004.
\end{thebibliography}

\end{document}